\section{Introducción}

\subsection{Contexto y Problemática}

La evaluación de calidad en call centers exige que los supervisores escuchen llamadas
completas para identificar momentos críticos de escalamiento emocional, lo cual es
ineficiente, limita la cobertura y provoca pérdida de oportunidades clave de
aprendizaje.

El análisis de sentimiento basado en texto pierde información valiosa contenida en la
forma en que se expresa algo. Características como el tono de voz, la velocidad del
habla y la tensión vocal revelan estados emocionales que el contenido textual no
captura, especialmente en casos de sarcasmo, ansiedad encubierta o frustración
contenida.

\subsection{Objetivos}

\textbf{Objetivo General:} Desarrollar un sistema automatizado para detectar picos
emocionales en las llamadas de un call center, basado en el análisis de características
paralingüísticas del habla.

\textbf{Objetivos Específicos:}
\begin{itemize}
    \item Extraer características acústicas relevantes.
    \item Evaluar métodos para calcular líneas base acústicas individualizadas.
    \item Comparar algoritmos de detección de anomalías para identificar picos
          emocionales.
    \item Desarrollar un sistema para detectar picos emocionales y clasificar la
          valencia emocional.
    \item Validar el sistema con llamadas reales y expertos.
\end{itemize}
